\documentclass{article}
\usepackage[utf8]{inputenc}
\usepackage{amsmath, amssymb}
\usepackage{graphicx}
\usepackage{hyperref}
\usepackage{geometry}
\geometry{a4paper, margin=1in}

\title{Evolution and
Selection of
Quantitative Traits}
\author{Unknown Authors}
\date{}

\begin{document}
\maketitle

\begin{abstract}

\end{abstract}

\section{Introduction}
Bruce Walsh
University of Arizona

Michael Lynch
Arizona State University

\begin{figure}[h]
\centering
\includegraphics[width=0.8\linewidth]{example-image}
\caption{Figure}
\end{figure}

Sinauer Associates

NEW YORK OXFORD
OXFORD UNIVERSITY PRESS

\section{Cover Photo Credits (top to bottom)}
• Lithophane leeae Walsh 2009 (Noctuid moth). Photo by Bruce Walsh.

• Chuck (Angora rabbit). Photo by Lee Fulmer.

• Maize diversity. Photo by Nick Kaczmar.

- Bony armor variation in Gasterosteus aculeatus (the threespine stickleback). Photo by Rowan Barrett.

\section{OXFORD}
UNIVERSITY PRESS

Great Clarendon Street, Oxford, OX2 6DP,

United Kingdom

Oxford University Press is a department of the University of Oxford. It furthers the University's objective of excellence in research, scholarship, and education by publishing worldwide. Oxford is a registered trade mark of Oxford University Press in the UK and in certain other countries.

© Bruce Walsh and Michael Lynch, 2018

The moral rights of the authors have been asserted

First Edition published in 2018

Impression: 1

All rights reserved. No part of this publication may be reproduced, stored in a retrieval system, or transmitted, in any form or by any means, without the prior permission in writing of Oxford University Press, or as expressly permitted by law, by licence or under terms agreed with the appropriate reprographics rights organization. Enquiries concerning reproduction outside the scope of the above should be sent to the Rights Department, Oxford University Press, at the address above

You must not circulate this work in any other form and you must impose this same condition on any acquirer

Published in the United States of America by Oxford University Press 198 Madison Avenue, New York, NY 10016, United States of America

British Library Cataloguing in Publication Data
Data available

Library of Congress Control Number: 2017054086

ISBN 978-0-19-883087-0

DOI: 10.1093/oso/9780198830870.001.0001

Printed in Great Britain by Bell & Bain Ltd., Glasgow

Links to third party websites are provided by Oxford in good faith and for information only. Oxford disclaims any responsibility for the materials contained in any third party website referenced in this work.

Seek simplicity and distrust it. Whitehead (1920)

No efforts of mine could avail to make the book easy reading.
Fisher (1930. The genetical theory of natural selection)

\section{Contents}
PREFACE xxvii
I. INTRODUCTION 1
1. CHANGES IN QUANTITATIVE TRAITS OVER TIME 3
A Brief History of the Study of the Evolution of Quantitative Traits 4
The fusion of population and quantitative genetics 4
The ongoing fusion of molecular and quantitative genetics 7
The common thread between breeding and evolution in natural populations 9
Detecting selection in natural populations 10
The Theoretical Foundations of Evolutionary Biology 11
The completeness of evolutionary theory 11
Nonadaptive hypotheses and our understanding of evolution 14
Overview and Pathways Through This Volume 16
Evolution at one and two loci 17
Drift and quantitative traits 18
Short-term response on a single character 18
Selection in structured populations 19
Population-genetic models of trait response 20
Measuring selection on traits 20
Appendices 21
Volume 3 21
Notation 22
II. EVOLUTION AT ONE AND TWO LOCI 23
2. NEUTRAL EVOLUTION IN ONE- AND TWO-LOCUS SYSTEMS 25
The Wright-Fisher Model 26
Loss of Heterozygosity by Random Genetic Drift 28
Probabilities and Times to Fixation or Loss 31
The age of a Neutral Allele 33
Allele-frequency Divergence Among Populations 35
Buri's Experiment 36
Higher-order Allele-frequency Moments 40
Linkage Disequilibrium 43
Mutation-drift Equilibrium 44
The Detailed Structure of Neutral Variation 48
The infinite-alleles model and the associated allele-frequency spectrum 48
The infinite-sites model and the associated site-frequency spectrum 50
The Genealogical Structure of a Population 52
Mutation-migration-drift Equilibrium 54
Quantifying population structure:  $ F_{ST} $ 54
Mutation-migration-drift equilibrium values of  $ F_{ST} $ 55
3. THE GENETIC EFFECTIVE SIZE OF A POPULATION 59
General Considerations 60
Monoecy 61
Dioecy 63
Age Structure 65
Variable Population Size 67

vii

viii

Partial Inbreeding ..... 68
Special Systems of Mating ..... 69
Population Subdivision ..... 72
Selection, Recombination, and Hitchhiking Effects ..... 74
Effects from selection at unlinked loci ..... 74
Selective sweeps and genetic draft ..... 75
Background selection ..... 77
4. THE NONADAPTIVE FORCES OF EVOLUTION ..... 81
Relative Power of Mutation and Genetic Drift ..... 82
Nucleotide diversity ..... 82
Number of segregating sites ..... 83
Alternative approaches ..... 84
Empirical observations ..... 85
Relative Power of Recombination and Genetic Drift ..... 86
Number of recombinational events in a sample of alleles ..... 87
Other approaches for narrow genomic intervals ..... 89
Large-scale analysis ..... 90
Empirical observations ..... 92
Effective Population Size ..... 94
Temporal change in allele frequencies ..... 94
Single-sample estimators ..... 96
Empirical observations ..... 97
Mutation Rate ..... 98
Divergence analysis ..... 98
Short-term enrichment ..... 100
Evolution of the mutation rate ..... 103
Recombination Rate ..... 106
Evolution of the recombination rate ..... 107
General Implications ..... 109
5. THE POPULATION GENETICS OF SELECTION ..... 113
Single-locus Selection: Two Alleles ..... 113
Viability selection ..... 114
Expected time for allele-frequency change ..... 116
Differential viability selection on the sexes ..... 117
Frequency-dependent selection ..... 118
Fertility/fecundity selection ..... 118
Sexual selection ..... 119
Wright's Formula ..... 119
Adaptive topographies and Wright's formula ..... 122
Single-locus Selection: Multiple Alleles ..... 122
Marginal fitnesses and average excesses ..... 122
Equilibrium frequencies with multiple alleles ..... 124
Internal, corner, and edge equilibria and basins of attraction ..... 125
Wright's formula with multiple alleles ..... 126
Changes in genotypic fitnesses,  $ W_{ij} $, when additional loci are under selection ..... 127
Selection on Two Loci ..... 128
Dynamics of gamete-frequency change ..... 129
Gametic equilibrium frequencies, linkage disequilibrium, and mean fitness ..... 130
Results for particular fitness models ..... 131
Phenotypic stabilizing selection and the maintenance of genetic variation ..... 132
Selection on a Quantitative Trait Locus ..... 135
Monogenic traits ..... 136
Many loci of small effect underlying the character ..... 136

ix

6. THEOREMS OF NATURAL SELECTION:
RESULTS OF PRICE, FISHER, AND ROBERTSON 145
Price's General Theorem of Selection 146
The life and times of George Price 146
Price's theorem,  $ R_{z} = \sigma(w_{i}, z_{i}) + E(w_{i} \overline{\delta}_{i}) $ 147
The Robertson-Price identity,  $ S = \sigma(w, z) $ 150
The breeder's equation,  $ R = h^{2}S $ 151
Fisher's Fundamental Theorem of Natural Selection 154
The classical interpretation of Fisher's fundamental theorem,  $ R_{w} = \sigma_{A}^{2}(w) $ 154
What did Fisher really mean? 157
Implications of Fisher's Theorem for Trait Variation 158
Traits correlated with fitness have lower heritabilities 160
Traits correlated with fitness have higher levels of both
additive and residual variance 162
Nonadditive genetic variance for traits under selection 164
Robertson's Secondary Theorem of Natural Selection 165
1968 version:  $ R_{z} = \sigma_{A}(w, z) $ 165
1966 version:  $ R_{z} = \sigma(w, A_{z}) $ 166
Accuracy of the secondary theorem 166
Connecting Robertson's results with those of Price, Fisher, and Lush 167
The Breeder's Equation Framed Within the Price Equation 168
Partial covariance and the spurious response to selection 168
Parent-offspring regressions before and after selection 171
Heywood's decomposition of response 172
7. INTERACTION OF SELECTION, MUTATION, AND DRIFT 175
Selection and Mutation at Single Loci 175
Selection and Drift at Single Loci 179
Probability of fixation under additive selection 180
Probability of fixation under arbitrary selection 185
Fixation of overdominant and underdominant alleles 187
Expected allele frequency in a particular generation 190
Joint Interaction of Selection, Drift, and Mutation 191
Haldane's Principle and the Mutation Load 194
Fixation Issues Involving Two Loci 197
The Hill-Robertson effect 197
Mutations with contextual effects 199
Stochastic tunneling 200
8. HITCHHIKING AND SELECTIVE SWEEPS 207
Sweeps: A Brief Overview 207
Hitchhiking, sweeps, and partial sweeps 207
Selection alters the coalescent structure at linked neutral sites 208
Hard vs. soft sweeps 211
The Behavior of a Neutral Locus Linked to a Selected Site 212
Allele-frequency change 212
Reduction in genetic diversity around a sweep 217
The Messer-Neher estimator of s 220
Recovery of variation following a sweep 220
Effects of sweeps on the variance in microsatellite copy number 221
The site-frequency spectrum 222
The pattern of linkage disequilibrium 224

Age of a sweep ..... 226
Partial sweeps ..... 226
Impacts of inbreeding ..... 227
Summary: Signatures of a hard sweep ..... 227
Soft Sweeps ..... 227
Signatures of a soft sweep ..... 229
Sweeps using standing variation ..... 230
How likely is a sweep using standing variation? ..... 231
“Hardening” and “softening” of a sweep ..... 235
Recurrent mutation of the favorable allele cannot be ignored ..... 236
Impact of geographic structure ..... 239
Polygenic Adaptation ..... 242
Genome-Wide Impact of Repeated Selection at Linked Sites ..... 243
Association between levels of variation and localized recombination rates ..... 244
Impact of recurrent hard sweeps on levels of neutral variation ..... 245
A few large or many small sweeps? ..... 247
Selective interference and the Hill-Robertson effect ..... 250
Background selection: Reduction in variation under low recombination or selfing ..... 251
Background selection vs. recurrent selective sweeps ..... 254
Sweeps, background selection, and substitution rates ..... 255
Sweeps, background selection, and codon usage bias ..... 257
A paradigm shift away from the neutral theory of molecular evolution? ..... 262
9. USING MOLECULAR DATA TO DETECT SELECTION:
   SIGNATURES FROM RECENT SINGLE EVENTS ..... 265
An Overview of Strategies Based on Segregating Variation ..... 266
Attempts to account for departures from the equilibrium model ..... 268
SNP ascertainment bias ..... 268
SNP polarity assignment errors ..... 269
Background selection as the more appropriate null? ..... 269
Structure of the remainder of this chapter ..... 269
Allele-Frequency Change in a Single Population ..... 269
Allele-frequency change over two sample points: The Waples adjusted test ..... 270
Allele-frequency change over a time series: The Fisher-Ford test ..... 272
Schaffer's linear trend test ..... 274
Scans and simulation-based approaches ..... 275
Birthdate selection mapping (BDSM) ..... 275
Divergence Between Populations: Two-population Comparisons ..... 277
Divergence Between Populations:  $ F_{ST} $-based Tests ..... 278
Outlier approaches ..... 279
Tests based on  $ F_{ST} $-generated branch-lengths ..... 280
The Lewontin-Krakauer test: Basics ..... 280
Second-generation Lewontin-Krakauer tests: Model-based adjustments ..... 282
Third-generation Lewontin-Krakauer tests: Correcting for population structure ..... 283
Allele-Frequency Correlations With Environmental Variables ..... 284
Joost's spatial analysis method (SAM) ..... 284
Accounting for population structure: Coop's Bayenv and Frichot's LFMM ..... 286
Changes in the Chromosomal Pattern of Neutral Variation ..... 289
Simple visual scans for changes in nucleotide and STR diversity ..... 289
Tests based on STR variation across populations ..... 291
Tests of sweeps using bottleneck models ..... 292
Tests of sweeps using genomic positional information: CLRT-GOF ..... 293
Tests of sweeps using genomic positional information: "SweepFinder" ..... 295
Tests of sweeps using genomic positional information: XP-CLR ..... 297
Ascertainment issues ..... 299

xi

Model fragility: Demography, mutation, recombination, and gene conversion ..... 299
Tests Based on Site-frequency Spectrum Statistics ..... 300
Summary statistics based on infinite-sites models ..... 300
Tajima's D test ..... 303
Achaz's Y and Y* tests ..... 305
Fu and Li's D* and F* tests ..... 305
Fay and Wu's H test ..... 306
Zeng et al.'s E test ..... 307
Adjusting the null to account for nonequilibrium populations ..... 307
Support via a preponderance of evidence ..... 309
Recombination makes site-frequency tests conservative ..... 311
Haplotype-based Tests ..... 311
Defining and inferring haplotypes ..... 312
Overview of haplotype-based tests ..... 312
The Ewens-Watterson test ..... 314
Other infinite-alleles tests: Conditioning on  $ \widehat{\theta} $ ..... 315
Other infinite-alleles tests: Conditioning on S ..... 316
Garud et al.'s H $ _{12} $ and H $ _{2} $ tests ..... 318
Recombination and infinite-alleles-based tests ..... 318
Pairwise disequilibrium tests: Kelly's Z $ _{nS} $ and Kim and Nielsen's  $ \omega_{max} $ ..... 319
Contrasting allele-frequency vs. intra-allelic estimates of haplotype age ..... 320
Long haplotype tests using within-population data:
rEHH, LDD, iHS, nS $ _{L} $, SDS, and DIND ..... 324
Long haplotype tests using between-population data:
XP-EHH, Rsb, rHH, rMHH, and  $ \chi_{MD} $ ..... 328
Summary: Tests based on haplotype/LD information ..... 330
Searches for Selection: Humans ..... 330
Recent and current selection in humans ..... 331
Balancing selection in humans ..... 333
Searches for Selection: Domesticated Organisms ..... 334
The process of domestication ..... 335
Finding domestication and improvement genes in crops ..... 336
Domestication and improvement genes in rice ..... 338
Domestication and improvement genes in maize ..... 338
Relative strengths of selection on domestication vs. improvement genes ..... 340
Silkmoths and flies ..... 341
Constraints on finding domestication and improvement genes through selective signals ..... 341
0. USING MOLECULAR DATA TO DETECT SELECTION:
SIGNATURES FROM MULTIPLE HISTORICAL EVENTS ..... 343
Brief Overview of Divergence-based Tests ..... 344
A history of selection alters the ratio of polymorphic to divergent sites ..... 344
A history of positive selection alters the ratio of silent- to replacement-site substitution rates ..... 345
Divergence-based tests are biased toward conservative sites ..... 346
What fraction of the genome is under functional constraints? ..... 346
The HKA and McDonald-Kreitman Tests ..... 347
The Hudson-Kreitman-Aguadé (HKA) test ..... 347
The McDonald-Kreitman (MK) test: Basics ..... 350
The McDonald-Kreitman test: Caveats ..... 356
Dominance in fitness and the MK test ..... 359
Fluctuating selection coefficients and MK tests ..... 359
Recombinational bias in extended MK tests ..... 360
Estimating Parameters of Adaptive Evolution ..... 360

xii

Estimating the fraction,  $ \alpha $, of substitutions that are adaptive ..... 360
How common are adaptive substitutions? ..... 365
Estimating the rate,  $ \lambda $, of adaptive substitutions ..... 367
The Sawyer-Hartl Poisson Random Field Model ..... 369
Basic structure of the model ..... 369
Bayesian extensions ..... 374
INSIGHT analysis of human transcription factors ..... 376
Phylogeny-based Divergence Tests ..... 377
The  $ K_{a} $ to  $ K_{s} $ ratio,  $ \omega $ ..... 378
Parsimony-based ancestral reconstruction tests ..... 379
Maximum-likelihood-based codon tests ..... 380
Bayesian estimators of sites under positive selection ..... 383
Connecting the Parameters of Adaptive Evolution ..... 384
The Search for Selection: Closing Comments ..... 385
Caution is in order when declaring positive selection ..... 386
Curbing our enthusiasm ..... 387
III. DRIFT AND QUANTITATIVE TRAITS ..... 389
11. CHANGES IN GENETIC VARIANCE INDUCED BY DRIFT ..... 391
Response of Within-population Genetic Variance to Drift ..... 391
Complete additivity ..... 391
The effects of dominance ..... 392
Quadratic components for inbred populations ..... 393
One- and two-locus identity coefficients ..... 395
Impact of drift under nonadditive variance ..... 397
The effects of epistatis ..... 399
Sampling error ..... 403
Empirical data ..... 405
Covariance Between Inbred Relatives ..... 411
REML estimates ..... 416
Empirical observations ..... 417
Drift-mutation Equilibrium ..... 419
Subdivided populations ..... 422
12. THE NEUTRAL DIVERGENCE OF QUANTITATIVE TRAITS ..... 425
Short-term divergence ..... 426
Sampling error ..... 427
Sample variance confidence intervals ..... 429
Empirical observations ..... 431
Estimating the among-group variance ..... 432
Lande's constant variance test,  $ F_{CV} $ ..... 433
Long-term Divergence ..... 434
Effectively neutral divergence and the estimation of rates of mutational variance ..... 436
Testing the Null Hypothesis of Neutral Phenotypic Divergence: Rate-based Tests ..... 438
Rate tests ..... 439
Lande's Brownian-motion model of neutral trait evolution ..... 439
Tests based on the Brownian-motion model ..... 441
Ornstein-Uhlenbeck models ..... 445
Divergence in morphological traits ..... 446
Time Series Divergence Tests ..... 446
Tests for departures from symmetric random walks ..... 447
Hunt's approach for comparing different models ..... 448
Population Structure Based-tests:  $ Q_{ST} $ versus  $ F_{ST} $ ..... 452
QST: Partitioning additive variance over populations ..... 452

xiii

$P_{ST}$: Approximating $Q_{ST}$ with phenotypic data ..... 454
Testing $Q_{ST}$ versus $F_{ST}$ ..... 455
Empirical data ..... 458
Closing comments: $Q_{ST}$, $F_{STQ}$, and linkage disequilibrium ..... 460
Trait-augmented Marker-based Approaches: Tests Using QTL Information ..... 462
Leveraging QTL studies ..... 462
Orr's QTLST and QTLST-EE sign tests ..... 462
Applications of QTL sign tests ..... 465
Trait-augmented Marker-based Approaches: Tests Using GWAS Information ..... 467
Approaches based on combining signals ..... 467
Tests based on tSDS scores ..... 467
The Berg and Coop $Q_{x}$ test ..... 468
Divergence in Gene Expression ..... 469
Level of gene expression as a quantitative trait ..... 469
Rate-based tests for neutrality in divergence of gene expression ..... 470
Largely neutral evolution of expression levels in primates? ..... 473
Transcriptional $Q_{ST}$, $tQ_{ST}$ ..... 474
Cis vs. trans, local vs. distant, and allele-specific expression (ASE) ..... 474
Applications of sign-based tests to expression data ..... 475
Evolution of expression levels: Drift, directional, or stabilizing selection? ..... 477
IV. SHORT-TERM RESPONSE ON A SINGLE CHARACTER ..... 479
13. SHORT-TERM CHANGES IN THE MEAN:
1. THE BREEDER'S EQUATION ..... 481
Single-generation Response: The Breeder's Equation ..... 482
The breeder's equation: A general approximation for response ..... 482
The importance of linearity ..... 482
Response is the change in mean breeding value ..... 483
Response under sex-dependent parent-offspring regressions ..... 484
The selection intensity, $\bar{i}$ ..... 486
The Robertson-Price identity, $S=\sigma(w,z)$ ..... 486
Correcting for reproductive differences: Effective selection differentials ..... 487
Expanding the Basic Breeder's Equation ..... 488
Accuracy ..... 489
Reducing environmental noise: Stratified mass selection ..... 493
Reducing environmental noise: Repeated-measures selection ..... 494
Adjustments for overlapping generations ..... 496
Maximizing response under the breeder's equation ..... 497
Maximizing the economic rate of response ..... 499
Mean- versus variance-standardized response ..... 499
BLUP Selection ..... 500
The Multivariate Breeder's Equation ..... 500
Response with two traits ..... 501
Accounting for phenotypic correlations: The selection gradient ..... 501
Accounting for genetic correlations: The Lande equation ..... 502
Selection gradients and mean population fitness ..... 503
Limitations of the Breeder's Equation ..... 504
14. SHORT-TERM CHANGES IN THE MEAN:
2. TRUNCATION AND THRESHOLD SELECTION ..... 507
Truncation Selection ..... 507
Selection intensities and differentials under truncation selection ..... 507
Correcting the Selection Intensity for Finite Sample Sizes ..... 509

Response in Discrete Traits: Binary Characters ..... 512
The threshold/liability model ..... 512
Direct selection on the threshold, T ..... 516
The logistic regression model for binary traits ..... 516
BLUP selection with binary data: Generalized linear mixed models ..... 520
Response in Discrete Traits: Poisson-distributed Characters ..... 522

15. SHORT-TERM CHANGES IN THE MEAN:
3. PERMANENT VERSUS TRANSIENT RESPONSE ..... 525
Why all the Focus on  $ h^{2} $? ..... 525
Genetic Sources of Transient Response ..... 526
Additive epistasis ..... 526
Dominance in autotetraploids ..... 529
Ancestral Regressions ..... 530
Response due to Environmental Correlations ..... 534
Selection in the Presence of Heritable Maternal Effects ..... 537
Decomposing maternal effects ..... 537
Selection response under Falconer's dilution model ..... 538
Separate direct and maternal traits ..... 543
Maternal selection vs. maternal inheritance ..... 545

16. SHORT-TERM CHANGES IN THE VARIANCE:
1. CHANGES IN THE ADDITIVE VARIANCE ..... 549
Changes in Variance due to Gametic-Phase Disequilibrium ..... 549
Changes in Variance Under the Infinitesimal Model ..... 552
Within- and among-family variance under the infinitesimal model ..... 554
Accounting for inbreeding and drift ..... 556
Changes in Variance Under Truncation Selection ..... 557
Changes in correlated characters ..... 560
Directional truncation selection: Theory ..... 561
Directional truncation selection: Experimental results ..... 563
Effects of epistasis: Does the Griffing effect overpower the Bulmer effect? ..... 563
Double truncation selection: Theory ..... 564
Double truncation selection: Experimental results ..... 565
Response under Normalizing Selection ..... 567
Selection with Assortative Mating ..... 569
Results using the infinitesimal model ..... 569
Assortative mating and enhanced response ..... 570
Disruptive selection, assortative mating, and reproductive isolation ..... 572

17. SHORT-TERM CHANGES IN THE VARIANCE:
2. CHANGES IN THE ENVIRONMENTAL VARIANCE ..... 573
Background: Heritable Variation in  $ \sigma_{E}^{2} $ ..... 573
Scales of environmental sensitivity ..... 573
Environmental vs. genetic canalization ..... 574
Evidence for heritable variation in the environmental variance ..... 575
Modeling Genetic Variation in  $ \sigma_{E}^{2} $ ..... 576
The multiplicative model ..... 577
The exponential model ..... 578
The additive model ..... 580
The heritability of the environmental variance,  $ h_{v}^{2} $ ..... 581
Selection on  $ \sigma_{E}^{2} $ ..... 583
Translating the response in  $ A_{v} $ into response in  $ \sigma_{E}^{2} $ ..... 583
Response from stabilizing selection on phenotypic value, z ..... 584

xv

Response from directional selection on z ..... 585
Direct selection on  $ \sigma_{E}^{2} $ using repeated records ..... 588

18. ANALYSIS OF SHORT-TERM SELECTION EXPERIMENTS:
1. LEAST-SQUARES APPROACHES ..... 591
Variance in Short-term Response ..... 591
Expected variance in response generated by drift ..... 592
Variance in predicted response versus variance in actual response ..... 595
Realized Heritabilities ..... 595
Estimators for several generations of selection ..... 596
Weighted least-squares estimates of realized heritability ..... 598
Standard errors of realized heritability estimates ..... 600
Power: Estimation of  $ h^{2} $ from relatives or selection response? ..... 601
Empirical versus predicted standard errors ..... 603
Realized heritability with rank data ..... 604
Infinitesimal-model corrections for disequilibrium ..... 604
Experimental Evaluation of the Breeder's Equation ..... 606
Most traits respond to selection ..... 606
Sheridan's analysis ..... 607
Realized heritabilities, selection intensity, and inbreeding ..... 608
Asymmetric selection response ..... 611
Reversed response ..... 615
Control Populations and Experimental Designs ..... 616
Basic theory of control populations ..... 616
Divergent selection designs ..... 618
Variance in response ..... 618
Control populations and variance in response ..... 619
Optimal Experimental Designs ..... 620
Nicholas' criterion ..... 621
Replicate lines ..... 621
Line-cross Analysis of Selection Experiments ..... 622
The simple additive model ..... 622
The Hammond-Gardner model ..... 623
Accounting for inbreeding depression and drift ..... 628

9. ANALYSIS OF SHORT-TERM SELECTION EXPERIMENTS:
2. MIXED-MODEL AND BAYESIAN APPROACHES ..... 631
Mixed-model vs. Least-squares Analysis ..... 631
BLUP selection ..... 632
Basics of Mixed-model Analysis ..... 633
REML estimation of unknown variance components ..... 636
REML often returns variance estimates unbiased by selection ..... 637
Animal-model Analysis of Selection Experiments ..... 638
The basic animal model ..... 638
Response is measured by change in mean breeding values ..... 639
Fixed effects alter heritabilities ..... 642
Model validation ..... 642
Separating genetic and environmental trends ..... 643
Validation that a trend is indeed genetic ..... 645
Replicate lines ..... 647
Estimating the additive variance at generation t ..... 647
The Relationship Matrix, A, Accounts for Drift and Disequilibrium ..... 648
Modifications of the Basic Animal Model ..... 651
Models with additional random effects ..... 652
Common family and maternal effects ..... 654

xvi

Treating certain breeding values as fixed effects ..... 655
Dominance ..... 656
Bayesian Mixed-model Analysis ..... 659
Introduction to Bayesian statistics ..... 659
Computing posteriors and marginals: MCMC and the Gibbs sampler ..... 662
Bayesian analysis of the animal model ..... 665
Application: Estimating selection response in pig litter-size components ..... 666
LS, MM, or Bayes? ..... 668
20. SELECTION RESPONSE IN NATURAL POPULATIONS ..... 671
Evolution in Natural Populations: What Is the Target of Selection? ..... 672
Direct and correlated responses ..... 672
Environmentally generated correlations between fitness and traits ..... 674
The Fisher-Price-Kirkpatrick-Arnold Model for Evolution of Breeding Date ..... 676
Modifying the Breeder's Equation for Natural Populations ..... 677
Complications in the absence of environmental change ..... 678
Additional complications from environmental change ..... 682
Is a Focal Trait the Direct Target of Selection? ..... 682
Robertson's theorem: Response prediction without regard to the target of selection ..... 682
Robertson consistency tests ..... 683
Rausher's consistency criteria ..... 684
Morrissey et al.'s consistency criteria ..... 686
The breeder's equation versus the secondary theorem ..... 687
Applying Mixed Models to Natural Populations: Basics ..... 688
Animal-model analysis in natural populations: Overview ..... 689
Obtaining the relationship matrix: Direct observation of the pedigree ..... 691
Obtaining the relationship matrix: Marker data ..... 691
Consequences of pedigree errors ..... 696
Model selection ..... 697
Applying Mixed Models to Natural Populations ..... 698
Consistency tests: Accuracy, reliability, and caveats when using PBVs ..... 698
Bivariate animal models: REML estimates of  $ \sigma(A_{z}, A_{w}), S_{A}, \text{and } S_{z} $ ..... 702
Next-generation analysis: Generalized mixed models ..... 703
Next-generation analysis: Modeling missing data ..... 704
Detecting genetic trends ..... 706
Causes of Apparent Failures of Response in Natural Populations ..... 709
Cryptic evolution: Genetic change masked by environmental change ..... 711
Antler size in red deer: The focal trait is not the target of selection ..... 713
Lower heritabilities in environments with stronger selection? ..... 713
Fitness tradeoffs and multivariate constraints ..... 714
V. SELECTION IN STRUCTURED POPULATIONS ..... 717
21. FAMILY-BASED SELECTION ..... 719
Introduction to Family-based Selection Schemes ..... 719
Overview of the different types of family-based selection ..... 720
Plant versus animal breeding ..... 720
Among- versus within-family selection ..... 722
Details of Family-based Selection Schemes ..... 722
Selection and recombination units ..... 722
Variations of the selection unit ..... 723
Variations of the recombination unit ..... 725
Theory of Expected Single-cycle Response ..... 727
Modifications of the breeder's equation for predicting family-based response ..... 727

The selection unit-offspring covariance,  $ \sigma(x, y) $ ..... 731
Variance of the selection unit,  $ \sigma_{x}^{2} $ ..... 733
Response for Particular Designs ..... 736
Overview of among- and within-family response ..... 736
Among-family selection ..... 737
Among-family selection: Which scheme is best? ..... 740
Within-family selection ..... 741
Realized heritabilities ..... 743
Accounting for epistasis ..... 744
Response with autotetraploids ..... 745
Efficiency of Family-based Vs. Individual Selection ..... 749
The relative accuracies of family-based vs. individual selection ..... 749
Comparing selection intensities: Finite size corrections ..... 753
Within-family selection has additional long-term advantages ..... 754
Response when Families are Replicated Over Environments ..... 755
Among-family variance under replication ..... 755
Ear-to-row selection ..... 758
Modified ear-to-row selection ..... 758
Selection on a Family Index ..... 761
Response to selection on a family index ..... 762
Lush's optimal index ..... 763
Correcting the selection intensity for correlated variables ..... 766
22. ASSOCIATIVE EFFECTS: COMPETITION, SOCIAL INTERACTIONS, GROUP AND KIN SELECTION ..... 769
Direct Versus Associative Effects ..... 770
Early models of competition ..... 771
Direct and associative effects ..... 771
Animal well-being and the improvement of the heritable social environment ..... 773
What do we mean by group? ..... 773
Trait- vs. variance component-based models ..... 774
The total breeding value (TBV) and  $ T^{2} $ ..... 775
 $ A_{s} $ as a function of group size ..... 778
Selection In the Presence of Associative Effects ..... 779
Individual selection: Theory ..... 780
Individual selection: Direct vs. social response ..... 782
Individual selection: Maternal effects ..... 783
Group selection: Theory ..... 784
Group selection: Direct vs. social response ..... 788
Group selection: Experimental evidence ..... 788
Incorporating Both Individual and Group Information ..... 789
Response on a weighted index ..... 789
Optimal response ..... 793
BLUP Estimation of Direct and Associative Effects ..... 796
Mixed-model estimation of direct and associative effects ..... 796
Muir's experiment: BLUP selection for quail weight ..... 800
Details: Environmental group effects and the covariance structure of e ..... 802
Details: Ignoring additive social values introduces bias ..... 804
Details: Identifiability of variance components ..... 805
Appropriate designs for estimating direct and associative effects ..... 806
Using kin groups: A quick-and-dirty way around associative effects? ..... 807
Associative Effects, Inclusive Fitness, and Fisher's Theorem ..... 807
Change in mean fitness when associative effects are present ..... 807
Inclusive fitness ..... 810
Bijma's theorem: Inclusive fitness and Fisher's fundamental theorem ..... 812

Hamilton's rule ..... 813
How general is Hamilton's rule? ..... 813
Queller's generalization of Hamilton's rule ..... 815
Group Selection, Kin Selection, and Associative Effects ..... 815
Kin, group, and multilevel selection ..... 815
Much ado about nothing? ..... 816
Group and kin selection: Models without trait associative effects ..... 817
Group and kin selection in the associative-effects framework ..... 821
Closing comments ..... 823
23. SELECTION UNDER INBREEDING ..... 825
Basic Issues in Selection Response Under Inbreeding ..... 826
Accounting for inbreeding depression ..... 826
Response under small amounts of inbreeding ..... 827
Using ancestral regressions to predict response ..... 828
The covariance between inbred relatives ..... 830
Limitations ..... 830
Family Selection with Inbreeding and Random Mating ..... 831
Family selection using inbred parents ..... 833
Progeny testing using inbred offspring ..... 835
S₁, S₂, and Sᵢ,ₙ family selection ..... 836
Other inbreeding-based family-selection schemes ..... 840
Cycles of inbreeding and outcrossing ..... 840
Individual Selection Under Pure Selfing ..... 842
Response under pure selfing ..... 843
Response when inbreeding pure-line crosses ..... 846
The Bulmer effect under selfing ..... 847
Family Selection Under Pure Selfing ..... 849
The covariance between relatives in a structured selfing population ..... 850
Response to family selection ..... 853
Within-family selection under selfing ..... 855
Combined selection ..... 856
Response Under Partial Selfing ..... 859
An approximate treatment using covariances ..... 859
A more careful treatment: Kelly's structured linear model ..... 861
Conclusions: An Incomplete Theory for Short-term Response ..... 864
The Evolution of Selfing Rates ..... 865
Automatic selection, inbreeding depression, and reproductive assurance ..... 865
The Lande-Schemske model: Theory ..... 867
The Lande-Schemske model: Data ..... 867
Baker's Law and the demographic advantages of selfing ..... 869
Group-level selection against selfing ..... 869
VI. POPULATION-GENETIC MODELS OF TRAIT RESPONSE ..... 873
24. THE INFINITESIMAL MODEL AND ITS EXTENSIONS ..... 875
The Infinitesimal: A Family of Models ..... 876
The Infinitesimal Genetics Model: Empirical Data ..... 877
Theoretical Implications of the Infinitesimal Genetics Model ..... 879
Selection does not change allele frequencies ..... 879
Accounting for dominance ..... 880
Disequilibrium ..... 881
Theoretical Implications of the Standard Infinitesimal Model ..... 881

What generates a Gaussian distribution within a family? ..... 881
Modifications of the Fisher-Bulmer infinitesimal model ..... 882
Gaussian Continuum-of-alleles Models ..... 883
Infinite-alleles and continuum-of-alleles models ..... 883
Drift ..... 884
Drift and a finite number of loci ..... 885
The effective number of loci,  $ n_{e} $ ..... 887
Dynamics:  $ \sigma_{a}^{2} $ and d change on different time scales ..... 888
How robust is the Gaussian continuum-of-alleles model? ..... 889
The Bulmer Effect Under Linkage ..... 890
An approximate treatment ..... 890
A more careful treatment ..... 892
Response Under Non-Gaussian Distributions ..... 895
Describing the genotypic distribution: Moments ..... 895
Describing the genotypic distribution: Cumulants and Gram-Charlier series ..... 898
Application: Departure from normality under truncation selection ..... 900
Short-term response ignoring linkage disequilibrium ..... 902
Gaussian versus rare-alleles approximations ..... 905
Short-term response ignoring allele-frequency change ..... 909
Effects of linkage ..... 911
Summary: Where Does All This Modeling Leave Us? ..... 911
25. LONG-TERM RESPONSE:
1. DETERMINISTIC ASPECTS ..... 913
Idealized Long-term Response in a Large Population ..... 913
Deterministic Single-locus Theory ..... 916
Expected contribution from a single locus ..... 916
Dudley's estimators of a, n, and  $ p_{0} $ ..... 917
Dynamics of allele-frequency change ..... 918
Major Genes Versus Polygenic Response: Theory ..... 922
Lande's model: Response with a major gene in an infinitesimal background ..... 922
Major Genes Versus Polygenic Response: Data ..... 927
Major genes appear to be important in response to anthropogenically induced selection ..... 927
What is the genetic architecture of response in long-term selection experiments? ..... 928
An Overview of Long-term Selection Experiments ..... 929
Estimating selection limits and half-lives ..... 930
General features of long-term selection experiments ..... 932
The nature of selection limits ..... 935
Increases in Variances and Accelerated Responses ..... 937
Rare alleles ..... 937
Major mutations ..... 938
Scale and environmental impacts on variances ..... 941
Linkage effects ..... 941
Epistasis ..... 943
Conflicts Between Natural and Artificial Selection ..... 944
Accumulation of lethals in selected lines ..... 946
Lerner's model of genetic homeostasis ..... 949
Artificial selection countered by natural stabilizing selection ..... 950
26. LONG-TERM RESPONSE:
2. FINITE POPULATION SIZE AND MUTATION ..... 953
The Population Genetics of Selection and Drift ..... 954
Fixation probabilities for alleles at a QTL ..... 955
Increased recombination rates following selection ..... 956
The Effect of Selection on Effective Population Size ..... 956

xx

The expected reduction in  $ N_{e} $ from directional selection ..... 957
Drift and Long-term Selection Response ..... 961
Basic theory ..... 961
Robertson's theory of selection limits ..... 962
Tests of Robertson's Theory of Selection Limits ..... 965
Weber's selection experiments on Drosophila flight speed ..... 967
The Effects of Linkage on the Selection Limit ..... 969
Optimal Selection Intensities for Maximizing Long-term Response ..... 972
Effects of Population Structure on Long-term Response ..... 974
Founder effects and population bottlenecks ..... 974
Population subdivision ..... 977
Within-family selection ..... 979
Asymptotic Response due to Mutational Input ..... 981
Results for the infinitesimal model ..... 981
Expected asymptotic response under more general conditions ..... 985
Additional models of mutational effects ..... 987
Optimizing the asymptotic selection response ..... 988
27. LONG-TERM RESPONSE:
3. ADAPTIVE WALKS ..... 991
Fisher's Model: The Adaptive Geometry of New Mutations ..... 991
Fisher-Kimura-Orr adaptive walks ..... 994
The cost of pleiotropy ..... 999
Adaptive walks under a moving optimum ..... 1001
Walks in Sequence Space: Maynard-Smith-Gillespie-Orr Adaptive Walks ..... 1002
SSWM models and the mutational landscape ..... 1003
Extreme-value theory (EVT) ..... 1003
Structure of adaptive walks under the SSWM model ..... 1005
The fitness distribution of beneficial alleles ..... 1009
Fisher's Geometry or EVT? ..... 1010
The Importance of History in Walks: Lessons From Lenksi's LTEE ..... 1012
28. MAINTENANCE OF QUANTITATIVE GENETIC VARIATION ..... 1015
Overview: The Maintenance of Variation ..... 1015
Maintaining genetic variation for quantitative traits ..... 1016
Mutation-drift Equilibrium ..... 1018
Mutational models and quantitative variation ..... 1018
Maintenance of Variation by Direct Selection ..... 1021
Fitness models of stabilizing selection ..... 1021
Stabilizing selection on a single trait ..... 1023
Stabilizing selection on multiple traits ..... 1024
Stabilizing selection countered by pleiotropic overdominance ..... 1026
Fluctuating and frequency-dependent stabilizing selection ..... 1027
Summary of direct-selection models ..... 1028
Neutral Traits With Pleiotropic Overdominance ..... 1028
Mutation-stabilizing Selection Balance: Basic Models ..... 1030
Latter-Bulmer diallelic models ..... 1031
Kimura-Lande-Fleming continuum-of-alleles models ..... 1033
Gaussian versus house-of-cards approximations for continuum-of-alleles models ..... 1035
Epistasis ..... 1039
Effects of linkage and mating systems ..... 1039
Spatial and temporal variation in the optimum ..... 1040
Summary: Implications of Gaussian versus HCA approximations ..... 1042
Mutation-stabilizing Selection Balance: Drift ..... 1050
Impact on equilibrium variances ..... 1050

xxi

Near neutrality at the underlying loci? ..... 1051
Mutation-stabilizing Selection Balance: Pleiotropy ..... 1052
Gaussian results ..... 1053
HCA results ..... 1055
Maintenance of Variation by Pleiotropic Deleterious Alleles ..... 1058
The Hill-Keightley pleiotropic side-effects model ..... 1058
Deleterious pleiotropy-stabilizing selection (joint-effects) models ..... 1065
How Well do the Models Fit the Data? ..... 1069
Strength of selection: Direct selection on a trait ..... 1070
Strength of selection: Persistence times of new mutants ..... 1073
Number of loci and mutation rates ..... 1074
What Does Genetic Architecture Tell Us? ..... 1075
Accelerated responses in artificial selection experiments ..... 1076
Kelly's test for rare recessives ..... 1077
Summary: What Forces Maintain Quantitative-genetic Variation? ..... 1077
VII. MEASURING SELECTION ON TRAITS ..... 1079
29. INDIVIDUAL FITNESS AND THE MEASUREMENT OF UNIVARIATE SELECTION ..... 1081
Episodes of Selection and the Assignment of Fitness ..... 1082
Fitness components ..... 1082
Assigning fitness components ..... 1083
Potential issues with assigning discrete fitness values ..... 1086
Assigning components of offspring fitness to their mothers ..... 1086
Concurrent selective episodes, reproductive timing, and individual fitness,  $ \lambda_{ind} $ ..... 1087
Sensitivities and elasticities of the elements of L ..... 1090
Variance in Individual Fitness ..... 1091
Partitioning I across episodes of selection ..... 1094
Correcting lifetime reproductive success for random offspring mortality ..... 1095
Caveats in using the opportunity for selection ..... 1095
Measuring Sexual Selection ..... 1097
Bateman's principles ..... 1098
Variance in mating success ..... 1099
The sexual selection, or Bateman, gradient ..... 1102
Describing Phenotypic Selection: Introductory Remarks ..... 1104
Fitness surfaces and landscapes ..... 1104
Describing Phenotypic Selection: Changes in Phenotypic Moments ..... 1106
Henshaw's distributional selection differential (DSD) ..... 1107
Directional selection ..... 1107
Quadratic selection ..... 1108
Under trait normality, gradients describe the local geometry of the fitness landscape ..... 1109
Under trait normality, gradients appear in selection response equations ..... 1110
Partitioning changes in means and variances into episodes of selection ..... 1110
Choice of the reference population: Independent partitioning ..... 1112
Standard errors for estimates of differentials and gradients ..... 1113
Describing Phenotypic Selection: Individual Fitness Surfaces ..... 1114
Linear and quadratic approximations of  $ W(z) $ ..... 1115
Lande-Arnold fitness regressions ..... 1117
The geometry of quadratic fitness functions ..... 1119
Hypothesis testing, approximate confidence intervals, and model validation ..... 1120
Power ..... 1122
Mean-standardized gradients and fitness elasticities ..... 1123
Moving Away From Quadratic Fitness Functions ..... 1125

xxii

Quadratic surfaces can be very misleading ..... 1125
Gaussian and exponential fitness functions ..... 1126
Semiparametric approaches: Schluter's cubic-spline estimate ..... 1128
Janzen-Stein and Morrissey-Sakrejda gradients: Calculating average and landscape gradients under general fitness functions ..... 1128
More Realistic Models of the Distribution of Fitness Components Generalized linear models for fitness components ..... 1129
Aster models: Modeling the distribution of total fitness ..... 1133
The Importance of Experimental Manipulation Performance and eco-evolutionary surfaces ..... 1134
0. MEASURING MULTIVARIATE SELECTION ..... 1139
Selection on Multivariate Phenotypes: Differentials and Gradients ..... 1139
Changes in the mean vector: The directional selection differential, S ..... 1140
The directional selection gradient vector, β ..... 1140
Directional gradients, fitness surface geometry, and selection response ..... 1143
Changes in the covariance matrix: The quadratic selection differential matrix, C ..... 1144
The quadratic selection gradient matrix, γ ..... 1146
Quadratic gradients, fitness surface geometry, and selection response ..... 1148
Fitness surface curvature and within-generation changes in variances and covariances ..... 1148
Multidimensional Quadratic Fitness Regressions ..... 1149
Estimation, hypothesis testing, and confidence intervals ..... 1150
Regression packages and coefficients of γ ..... 1151
Geometric aspects ..... 1152
A brief digression: Orthonormal and diagonalized matrices ..... 1153
Canonical transformation of γ ..... 1154
Are traits based on canonical axes meaningful? ..... 1157
Strength of selection: γ_{ii} versus λ ..... 1158
Significance and confidence regions for a stationary point ..... 1159
Using Aster models to estimate fitness surfaces ..... 1159
Multivariate Semiparametric Fitness Surface Estimation ..... 1160
Projection-pursuit regression and thin-plate splines ..... 1162
Gradients for general fitness surfaces ..... 1163
Calsbeek's tensor approach for detecting variation in fitness surfaces ..... 1163
Model selection ..... 1164
The Strength and Pattern of Selection in Natural Populations ..... 1165
Meta-analysis ..... 1165
Kingsolver's analysis ..... 1166
Directional selection: Strong or weak? ..... 1168
Total versus direct selection, tradeoffs, and temporal variation ..... 1171
Directional selection on body size and Cope's law ..... 1173
Quadratic selection: Strong or weak? ..... 1173
Where is all the stabilizing selection? ..... 1175
A plea to fully publish ..... 1175
Unmeasured Characters and Other Biological Caveats ..... 1175
Path Analysis and Fitness Estimation ..... 1177
Regressions versus path analysis ..... 1179
Morrissey's extended selection gradient vector, η ..... 1181
Selection on latent variables ..... 1183
Elasticity path diagrams ..... 1184
Levels of Selection ..... 1186
Contextual analysis ..... 1186
Selection can be antagonistic across levels ..... 1187
Group selection is likely density-dependent ..... 1188
Selection differentials can be misleading in levels of selection ..... 1188

Hard, soft, and group selection: A contextual analysis viewpoint ..... 1190
Early survival: Offspring or maternal fitness component? ..... 1192
VIII. APPENDICES ..... 1197
A1. DIFFUSION THEORY ..... 1199
Foundations of Diffusion Theory ..... 1199
The infinitesimal mean,  $ m(x) $, and variance,  $ v(x) $ ..... 1199
The Kolmogorov forward equation ..... 1200
Boundary behavior of a diffusion ..... 1202
Derivation of the Kolmogorov forward equation ..... 1202
Stationary distributions ..... 1203
The Kolmogorov backward equation ..... 1205
Diffusion Applications in Population Genetics ..... 1205
Probability of fixation ..... 1206
Time to fixation ..... 1208
Expectations of more general functions ..... 1211
Applications in Quantitative Genetics ..... 1212
Brownian-motion models ..... 1212
Ornstein-Uhlenbeck models ..... 1215
A2. INTRODUCTION TO BAYESIAN ANALYSIS ..... 1217
Why are Bayesian Methods Becoming More Popular? ..... 1217
Bayes' Theorem ..... 1218
From Likelihood to Bayesian Analysis ..... 1220
Marginal posterior distributions ..... 1221
Summarizing the posterior distribution ..... 1221
Highest density regions (HDRs) ..... 1222
Bayes factors and hypothesis testing ..... 1222
The Choice of a Prior ..... 1224
Diffuse priors ..... 1224
The Jeffreys prior ..... 1225
Posterior Distributions Under Normality Assumptions ..... 1227
Gaussian likelihood with known variance and unknown mean ..... 1227
Gamma,  $ \chi^{2} $, inverse-gamma, and  $ \chi^{-2} $ distributions ..... 1228
Gaussian likelihood with unknown variance: Scaled inverse- $ \chi^{2} $ priors ..... 1231
Student's t distribution ..... 1232
General Gaussian likelihood: Unknown mean and variance ..... 1232
Conjugate Priors ..... 1233
The beta and Dirichlet distributions ..... 1233
Wishart and inverse-Wishart distributions ..... 1235
A3. MARKOV CHAIN MONTE CARLO AND GIBBS SAMPLING ..... 1237
Monte Carlo Integration ..... 1237
Importance sampling ..... 1239
Introduction to Markov Chains ..... 1240
The Metropolis-Hastings Algorithm ..... 1243
Burning-in the sampler ..... 1246
Simulated annealing ..... 1247
Choosing a jumping (proposal) distribution ..... 1247
Autocorrelation and sample size inflation ..... 1248
The Monte Carlo variance of an MCMC-based estimate ..... 1250
Convergence Diagnostics ..... 1251
Visual analysis ..... 1251

More formal approaches ..... 1252
Practical MCMC: How many chains and how long should they run? ..... 1253
The Gibbs Sampler ..... 1253
Using the Gibbs sampler to approximate marginal distributions ..... 1255
Rejection Sampling and Approximate Bayesian Computation (ABC) ..... 1256
A4. MULTIPLE COMPARISONS: BONFERRONI CORRECTIONS, FALSE-DISCOVERY RATES, AND META-ANALYSIS ..... 1259
Combining p Values Over Independent Tests ..... 1260
Fisher's  $ \chi^{2} $ method ..... 1260
Stouffer's Z score ..... 1261
Bonferroni Corrections and Their Extensions ..... 1262
Standard Bonferroni corrections ..... 1262
Sequential Bonferroni corrections ..... 1263
Holm's method ..... 1263
Simes-Hochberg method ..... 1264
Hommel's method ..... 1264
Cheverud's method and other approaches for dealing with dependence ..... 1265
Detecting an Excess Number of Significant Tests ..... 1266
How many false positives? ..... 1266
Schweder-Spjötvoll plots ..... 1267
Estimating  $ n_{0} $: Subsampling from a uniform distribution ..... 1267
Estimating  $ n_{0} $: Mixture models ..... 1270
FDR: The False-discovery Rate ..... 1271
Morton's posterior error rate (PER) and the FDR ..... 1272
A technical aside: Different definitions of false-discovery rate ..... 1275
The Benjamini-Hochberg FDR estimator ..... 1276
A (slightly more) formal derivation of the estimated FDR ..... 1277
Storey's q value ..... 1278
Closing caveats in using the FDR ..... 1279
Formal Meta-analysis ..... 1279
Informal, or narrative, meta-analysis ..... 1279
Standardizing effect sizes ..... 1280
Fixed-effects, random-effects, and mixed-model meta-analysis ..... 1282
Publication and other sources in bias ..... 1284
Bias when estimating magnitudes ..... 1287
A5. THE GEOMETRY OF VECTORS AND MATRICES: EIGENVALUES AND EIGENVECTORS ..... 1291
The Geometry of Vectors and Matrices ..... 1291
Comparing vectors: Lengths and angles ..... 1291
Matrices describe vector transformations ..... 1293
Orthonormal matrices: Rigid rotations ..... 1294
Eigenvalues and eigenvectors ..... 1294
Properties of Symmetric Matrices ..... 1297
Correlations can be removed by a matrix transformation ..... 1299
Simultaneous diagonalization ..... 1301
Canonical Axes of Quadratic Forms ..... 1302
Implications for the multivariate normal distribution ..... 1304
Principal components of the variance-covariance matrix ..... 1305
Testing for Multivariate Normality ..... 1307
Graphical tests: Chi-square plots ..... 1307
Mardia's test: Multivariate skewness and kurtosis ..... 1309

A6. DERIVATIVES OF VECTORS AND VECTOR-VALUED FUNCTIONS 1311
Derivatives of Vectors and Vector-valued Functions 1311
The Hessian Matrix, Local Maxima/Minima, and Multidimensional Taylor Series 1315
Optimization Under Constraints 1317
LITERATURE CITED 1319
AUTHOR INDEX 1411
ORGANISM AND TRAIT INDEX 1427
SUBJECT INDEX 1437

One hundred years ago, Sir Ronald Aylmer Fisher (1918), in a paper that was rejected by the Royal Society of London, laid the foundations for the field that has come to be known as quantitative genetics. This paper, and subsequent work by Fisher and others (often in the search of biologically motivated problems), also gave birth to much of modern statistics. The past century has seen quantitative genetics flourish, and indeed, the field has never been stronger or more vibrant as it moves seamlessly into the post-genomics era with an ever-broadening reach. If Fisher were alive today, we believe that he would be fascinated, and somewhat pleased, both at how much the field has grown, and also at how much it remains very familiar to him, even after all these years.

This text is the second part (out of three) of our treatment of the modern field of quantitative genetics from its humble, but path-breaking, beginnings 100 years ago. Some 20 years ago, we remarked in the introduction to the first part of our treatment (Lynch and Walsh 1998) that “today’s quantitative genetics is not the science that it was 25 (or even 10) years ago,” a statement that is even truer today. Our further comment that “the current machinery of quantitative genetics stands waiting (and its practitioners willing) to incorporate the fine genetic details of complex traits being elucidated by molecular and developmental biologists” has certainly been proven true over the past two decades. Modern quantitative genetics is the glue that connects many disciplines, especially given its ability to model uncertainty while fully accounting for known genetic and genomic features.

The first part of our treatment (Volume 1) dealt with genetics, which is the underpinning of complex traits. While we also remark on some recent developments in this area in this volume, our main focus is on the evolution of such traits, with applied evolution (plant and animal breeding) as an important application. To address these concerns, we have attempted to merge modern population-genetics theory with quantitative genetics and genomics. As a result, some of our treatment is more technical than was presented in Volume 1, as it provides a much fuller development of mathematical models of evolution. We have strived to present a holistic treatment, showing the interplay between theory and data, with extensive discussions on statistical issues relating to the estimation of the biologically relevant parameters for these models. We view quantitative genetics as the bridge between complex mathematical models of trait evolution and real-world data, and have tried to frame our treatment as such.

As with our first volume, a key goal in producing this volume is to increase communication among the various disciplines that make up the very diverse field of quantitative genetics. Twenty years ago, this included plant and animal breeders, evolutionary biologists, human geneticists, and statisticians. This list has only grown with the advent of genomics, which attracts an ever-widening array of those who use quantitative genetics (many of whom may not even be aware of it).

Given a project of this size and scope, and despite our own best efforts, excellent copy editors, and wonderful colleagues, it is likely that errors have escaped scrutiny. Hopefully, most will be obvious and trivial, but the posterior for this process often has a long tail, and, therefore, the potential for some interesting outliers. Please report any errors you find to BW (jbwalsh@email.arizona.edu). Likewise, we have a facebook group, Quantitative Genetics Book, whose file section will contain updated pdf files for current errata, as well as a forum for discussing issues related to the book. Finally, the final part of our treatment, which deals largely with multivariate issues, will be forthcoming in a more timely manner than this volume!

xxvii

\section{Acknowledgments}
As numerous seasoned readers will know, various components of this book have been in the works for over three decades. For a project of this magnitude, there is a host of people to thank, and we have undoubtedly forgotten to credit individuals who provided us with critical feedback in the earliest years. We profusely apologize for any such slights and blame our slowly fading memories.

As we moved into the field of quantitative genetics in earnest in the 1980s, critical foundational support in the development of our ideas was provided by our colleagues and teachers Stevan Arnold, Reinhard Burger, James Crow, Joe Felsenstein, Wilfried Gabriel, Daniel Gianola, Thomas Hansen, William Hill, Elizabeth Housworth, Alex Kondrashov, Russell Lande, Bill Muir, Tom Nagylaki, Tim Prout, Monty Slaktin, Michael Turelli, and Michael Wade.

We have also greatly benefited from all past and current members of our labs for their comments over the years. These especially include Ph.D. students: Matthew Ackerman, Desiree Allen, Stephan Baehr, David Butcher, Chi Chun Chen, Hong-Wen Deng, Suzanne Estes, Allan Force, Jean-François Pierre Gout, Kyle Hagner, Parul Johri, Vaishali Katju, Weiyi Li, Timothy Licknack, Hongan Long, Samuel Miller, Kendall Morgan, Angela Omilian, Michael Pfrender, Taylor Raborn, Sarah Schaack, Ryan Stikeleather, and Ken Spitze; and post-doctoral associates: Charles Baer, Melania Cristescu, Dee Denver, Thomas Doak, Jeffrey Dudycha, Suzanne Edmands, Jean Francois Gout, Kevin Higgins, David Houle, Sibel Kucukyildirim, Niles Lehman, Hong-An Long, Takahiro Maruki, Martin O'Hely, Susan Ratner, Barrie Robison, Stewart Schultz, Douglas Scofield, and John Willis.

BW also wishes to thank the thousands of students who have attended his short courses on various aspects of quantitative genetics over the past two decades, during which he has taught over 5000 students from roughly 60 countries at various locations around the world (25 countries at last count). This international interaction and feedback have greatly sharpened our ability to present the material in this book to a truly universal audience.

We have been exceptionally blessed with outstanding feedback throughout this project from a long list of colleagues, who directly engaged with the manuscript. Sally Otto and Michael Morrissey read whole parts of the manuscript, while critical comments on specific chapters were provided by Steve Arnold, Nick Barton, Piter Bijma, Mark Kirkpatrick, Bill Muir, Shinichi Nakagawa, Guilherme Rosa, Ruth Shaw, John Storey, and Michael Turelli.

We especially want to acknowledge Bill Hill, who has been a major force throughout this project, and in many ways could be considered its third author. Bill read all of both Volume 1 (Lynch and Walsh) and the current volume, offering very detailed and critical comments, and saving us from much embarrassment. Our goal throughout the project was to present the field with the breadth and depth that Bill carries around in his head. He also kept us going by constantly ending his emails with the reminder to “keep writing!” Hence, we are very pleased to codenicate this volume to Bill.

With respect to the production end, we have greatly enjoyed working with the highly professional staff at Sinauer Associates through the bulk of the editing process, and then with their new colleagues at Oxford during the final phase. Copy editor Nicole Balant and production editor Martha Lorantos faced the thankless task of trying to catch typos in a 1400-page technical manuscript where the lead author is rather dyslexic.

Elizabeth Morales and Ann Chiara did a wonderful job with the illustrations and patiently dealt with us, as we frequently changed our minds on how to best illustrate various concepts, while Michele Beckta dealt with copyright issues regarding the figures. Last, but certainty not least, production manager Chris Small thoughtfully guided us through the typesetting of both this and the previous volume (both of which were typeset by BW using Textures). Any less-than-professional appearance of aspects of this volume are there despite the best efforts of Chris!

We save very special thanks for our publisher and friend, Andy Sinauer. Andy initially started this project in the mid-1980s when we (BW and ML) discovered (via Andy) that

we were both working on separate books on the same subject and he logically queried whether it might be better if we worked together. Indeed it was! Andy has shown incredible patience over the past three decades on the pace of this project, and we hope that he will be pleased with the final outcome, despite its long gestation. We are thus extremely pleased to codenicate this volume to Andy for his long and outstanding service as one of the preeminent publishers of biological science textbooks. Upon his retirement, he has left a gap that will be hard to fill, and his impact on a number of fields, particularly evolutionary biology, is hard to overstate.

In closing, ML is especially grateful to the administration and colleagues at the University of Illinois, his place of first employment, for providing him with a safe environment for making an early (and risky) career transition from limnology to evolutionary genetics. And we both thank NSF and NIH for continuing support for research related to the subject matter of the book.

Finally, and especially, we wish to thank our incredible wives, Emília Martins (ML) and Lee Fulmer (BW), for their enduring patience during this long project.

Bruce Walsh, Tucson

Michael Lynch, Tempe

March 2018

\end{document}